%%%%%%%%%%%%%%%%%%%%%%%%%%%%%%%%%%%%%%%%%%%%%%%%%%%%%%%%%%%%%%%%%%%%%%%%%%%%%%%%%%%%%%%%%%%%%%%%%%%%%%%%%%%%
%%%%%%%%%%%%%%%%%%%%%%%%%%%%%%%%%% CONCLUSION   %%%%%%%%%%%%%%%%%%%%%%%%%%%%%%%%%%%%%%5%%%%%%%%%%%%%%%%%%%%%
%%%%%%%%%%%%%%%%%%%%%%%%%%%%%%%%%%%%%%%%%%%%%%%%%%%%%%%%%%%%%%%%%%%%%%%%%%%%%%%%%%%%%%%%%%%%%%%%%%%%%%%%%%%%

%Notably, the data coverage is more extensive for subsidiary-protected  individuals compared to Geneva Convention (GC) refugees. For the years with most asylum-seekers arriving in Austria, the coverage is especially commendable, reaching 67\% for GC refugees and 96\% for subsidiary-protected  individuals in the year 2015. The arrivals years with the highest number of asylum applications and the highest shares of coverage in our data, are also the main arrival years for our analysis sample. This reduces the possible selection bias. 


\FloatBarrier
\section{\label{sec:conclusion}Conclusion}

This study examines the interplay between local labor market conditions and welfare generosity on the labor market integration and internal mobility of refugees. Utilizing detailed register data, we analyze how refugees respond to the conditions they encounter when their choice set expands upon receiving protection.
Many initially dispersed refugees relocate as soon as they are allowed to do so. Their decision to relocate is influenced by both labor demand and welfare generosity at their initial placement location.

We show that labor demand, particularly in occupations with low entry barriers, immediately and significantly increases labor force participation among men. However, this effect fades after 2.5 years. In turn, tighter labor markets raise the likelihood of staying in the assigned state permanently.
Higher benefit levels decrease the likelihood of relocation but temporarily lower men’s employment levels and more persistently affect women’s employment, although the effect sizes are relatively small.

Lastly, when examining the interaction between labor market tightness and benefit levels, we find that changes in welfare benefits have a stronger effect on employment rates when local labor markets are tight and are negligible otherwise, consistent with similar findings for Denmark \citep{dustmann2024}.

By examining both relocation and employment, we show that tighter local labor markets encourage refugees to remain in their state and secure employment rather than relocate to another state without employment. Higher benefit levels have a minor negative effect on the probability of securing employment in the initial state but increase the likelihood of remaining without employment instead of moving and being unemployed elsewhere. Overall, our findings suggest that mobility responses to initial labor-market conditions and welfare benefit levels are more reflective of short-term needs than long-term optimization regarding labor market prospects, resembling the findings of work-first policies \citep{arendt_labor_2022, arendt_trade-offs_2023}.

From a policy perspective, our findings show that supporting refugees in finding employment at the moment of, or even before, receiving protection may prevent their migration to other regions. Additionally, they highlight the importance of considering the labor market conditions in the districts to which refugees are allocated. At the same time, changes in benefits will only affect labor market integration in strong labor markets. Thus, welfare reforms aimed at incentivizing employment among refugees will be effective only in regions with strong labor markets.

\begin{comment}
Our data revealed a positive effect of the $IVUR$ and refugee employment. A shift of two quartiles in the $IVUR$ leads to a 4.2 percentage point increase in the likelihood of employment by the 12th month, a substantial 37\% rise relative to the baseline employment rate in the same period. On the contrary, a similar two quartile shift in the $IPBL$ results in a 2.3 percentage point decrease in the employment probability, a 20.3\% reduction from the baseline.

Monthly Real Earnings: An upward shift in the $IVUR$ from the first to the third quartile corresponds to an increase of €84.6 in monthly real earnings, a notable 41\% rise relative to the mean in month 12. In contrast, a similar shift in the $IPBL$ leads to a reduction of €50.01 in earnings, a 24\% decline from the month 12 mean. Furthermore, every increase of €100 in potential benefits results in an €11 reduction in monthly real earnings.
In simpler terms, regions with more job openings relative to the number of unemployed individuals offer better employment opportunities for refugees. This observation underscores the pivotal role of immediate job availability in early labor market integration. Conversely, higher potential benefit levels tend to discourage refugees from entering the labor market promptly. While these benefits are crucial for the well-being of refugees, they might also serve as a deterrent to seeking employment. This highlights the trade-offs refugees face between labor market participation and accessing socio-economic support. However, our analysis does not conclusively determine whether rapid integration into the labor market might compel refugees to accept jobs that do not align with their qualification profiles, potentially impeding their long-term trajectories.

Interestingly, the effects of both the $IVUR$ and $IPBL$ on economic outcomes were found to be short-lived, losing significance within 24 months after refugees gained labor market access. This suggests that while initial conditions play a role, their impact diminishes over time.

Expanding on mobility outcomes, alterations in the $IPBL$ and $IVUR$ significantly influenced the likelihood of refugees staying in their initially assigned states. Notably, a 0.456 change in the $IPBL$ resulted in a 4.2 percentage point increase and a 0.211 change in the $IVUR$ resulted in a 4.1 percentage point increase, with both effects being comparatively persistent over time. Furthermore, increased potential benefits reduced the likelihood of refugees relocating to Vienna, signaling the role of benefits in influencing mobility patterns.
When observing social support outcomes, there was no significant effect from either the $IVUR$ or $IPBL$ on the utilization of public employment service (PES) services. However, a shift in the $IVUR$ from the first to the third quartile reduced the likelihood of refugees receiving welfare benefits by 4.9 percentage points, reinforcing the inverse relationship between employment and reliance on welfare benefits.
From a policy perspective, our findings highlight the importance of providing adequate job opportunities in regions with significant refugee placements. At the same time, it's vital to structure benefits in a way that supports refugees without discouraging labor market participation. Additionally, ensuring robust social support structures is an important aspect when it comes to social integration, potentially reducing the need to relocate to other parts of a country due to financial restraints. 
In conclusion, our research sheds light on the multifaceted aspects of refugee labor market integration in the Austrian context. The balance between job opportunities, socio-economic benefits, mobility, and social support, along with the time-bound nature of their effects, offers important insights for policymakers and stakeholders involved in refugee integration.
\end{comment}


